New technologies have made it easier than ever for individuals to credibly document granular details of their behavior. Firms can elicit this data to offer tailored products and personalized pricing, which may in turn benefit consumers and give rise to voluntary data sharing.

The potential benefits of behavioral data are especially compelling in the auto insurance industry. Granular driving data allows insurers to identify drivers with lower accident risk more precisely than algorithms based on public data, which pool consumers based on demographic features and sparse accident records alone \parencite{Cohen2007}. This can benefit both the insurer, who can more easily target safe drivers, and drivers, who may receive price discounts by demonstrating safe driving behavior.

However, consumers must volunteer to share their data in order for these benefits to manifest. Those who qualify for safe-driving discounts may not share their data due to fear of reclassification risk or other concerns \parencite{Handel2015}. Furthermore, firms typically retain property rights over the data that they collect. This incentivizes more data collection, but it also generates a competitive advantage that may lead to higher markups in the future \parencite{jin2021big}.

In this paper, we develop an empirical framework to examine the trade-offs for consumer welfare and firm profit that result from consumer data sharing in the context of an auto-insurance {monitoring} program  (``pay-how-you-drive'') in the United States. New customers are invited to plug a simple device into their cars, which tracks and reports their driving behavior for up to six months (\Cref{fig:eg_device}). In exchange, the insurer uses the data to better assess accident risk and adjust future premiums. Unlike traditional pricing factors such as age or accident records, monitoring data is not available to other firms. In 2017, insurers serving over 60\% of the \$267 billion U.S. auto insurance industry offered monitoring options.\footnote{\label{NAIC}2017 annual report of the National Association of Insurance Commissioners. Our research window is from 2012 to 2016.} Similar programs have been introduced in other industries, such as life insurance and consumer lending (Figure \ref{fig:eg_dtcd}).\footnote{The Vitality program from life insurer John Hancock tracks and rewards exercise and health-related behaviors. Ant Financial incentivizes users to conduct more personal finance transactions in exchange for borrowing discounts.} Since our data period, most major insurers have adopted smartphone-based monitoring programs; we discuss how the resulting changes in the competitive landscape affect the interpretation of our findings in Section~\ref{sec:related_lit_and_conclusion}. However, despite their growing relevance, empirical evidence on the effects of consumer data sharing mechanisms and their impact on welfare across different sides of the market is sparse.

We construct a novel dataset that combines proprietary individual-level panel data from a major U.S. auto insurer (hereinafter referred to as ``the firm'') and price menus offered by its competitors. Our panel covers 23 states and spans 2012 to 2016, during which the firm's monitoring program---the first major program in the industry---was introduced. For each consumer in our panel, we observe demographic characteristics, price menus from top competing insurers, insurance contracts purchased, and realized insurance claims. For each consumer who opts into the monitoring program, we observe a monitoring score reflecting her safe driving performance and the corresponding premium adjustments. Taken together, our analysis uses a panel dataset of over 1 million consumers and 50 million insurance quotes.

We first establish two key facts about the monitoring program as it is observed in our data. First, the monitoring program induces safer driving. Monitoring only occurs in the first period of insurance for new customers who opt in. Using a difference-in-differences estimator to capture within-consumer across-period variation in insurance claims, we find a 30\% \textit{moral hazard} effect: drivers who opt in to monitoring become safer on average. However, this behavioral change accounts for only 61\% of the risk differences between consumers in the monitored and unmonitored groups. Even after the monitoring period, the monitored group remains safer on average, and individual monitoring scores remain highly predictive of risk. This suggests that monitoring uncovers persistent, previously unobserved heterogeneity in risk, leading to \textit{advantageous selection} into the program.

These reduced-form analyses allow us to separately identify the effects of moral hazard and selection, extending a rich empirical literature that has thus far relied on correlations between plan choice and claims \parencite{Chiappori2000, Cohen2007, jeziorski2019skimming}. However,
quantifying the welfare impact of the monitoring program, or of prospective data regulations, requires a structural understanding of how heterogeneity in consumer risk and preferences shapes insurance demand, selection patterns, and firm pricing in equilibrium \parencite{Einav2010beyond}.
To do this, we develop an equilibrium model of insurance and monitoring.

On the demand side, we estimate a dynamic model that captures complex correlations between consumers' choice of insurance and monitoring opt-in, and their accident risk.
Consumers have heterogeneous private risk types and preferences---risk aversion, inertia costs, and disutility from being monitored---and choose an insurance plan from a personalized price menu with options offered by the firm and its competitors. When consumers engage with the firm for the first time, they can opt in to monitoring and receive an upfront discount. At the end of each period, accidents arrive based on consumers' risk types. In addition, if a consumer is monitored, a score is realized from a noisy distribution that is correlated with their risk type. In expectation, better scores lead to higher discounts in future periods. Consumers are thus incentivized to drive more safely when monitored.

On the supply side, we model firm pricing before and after the monitoring period to reflect the trade-off between eliciting more consumer data and profiting from this data in a competitive market. By offering a high upfront discount, the firm can encourage more consumers to opt in to monitoring. In order for this to be profitable, however, the firm must extract some of the surplus generated from monitoring in renewal periods. In this sense, the amount of monitoring data that is created (due to equilibrium monitoring participation) is endogenous to the firm's dynamic pricing strategy. In addition, the firm's ability to profitably ``invest-and-harvest" through the monitoring program is tempered by competition as it must attract and retain consumers in the first place.

Taken together, our model captures several forces that drive consumers' monitoring opt-in choice. First, consumers anticipate that their risk will be lower during the monitoring period if they opt in. Some may also expect higher future discounts and thus stand to gain more from monitoring. But the expected benefit of such future discounts is moderated by elevated reclassification risk due to noise in the monitoring process. Finally, consumers may incur various unobserved costs due to monitoring, including privacy loss, increased effort while driving, and additional decision-making. We quantify the combined effect of these costs with a heterogeneous monitoring disutility term.

Our model fits the data closely across key moments---claims, monitoring scores, pricing factors, coverage shares, monitoring opt-in, renewal attrition, and selection patterns.
The fit is stable across major product-menu changes, including the introduction of monitoring and a state-level mandatory-minimum increase.
Further, estimates from the pricing and score models demonstrate that monitoring sharpens risk rating: both premiums and scores increase with claims risk, leading the monitored group to face a steeper price-risk gradient and, consequently, to exhibit advantageous selection into monitoring.
The average consumer has modest risk aversion ($1.43\times10^{-5}$), but faces sizable switching frictions: \$333 per period for switching firms, or \$113 to opt in to monitoring.

Our counterfactual simulations indicate that both consumers and the firm benefit from the status quo compared to a scenario in which no monitoring program exists: total annual surplus increases by \$9.79 per capita per year even though the monitoring firm only captures 11.85\% of the market.
But monitoring take-up is low---due, in part, to choice frictions, reclassification risk, and competition from cheap plans offered by other insurers. The optimal pricing of monitoring, all else equal, would entail significantly more aggressive prices in the initial period---a 9\% surcharge and 77\% discount in the unmonitored and monitored pools, respectively---and much higher rent capture post monitoring among monitored drivers.
Due to the large surplus generated by the monitoring program and the firm's natural incentive to ``invest'' in it, counterfactual policies---such as enforcing a data portability regulation\footnote{This is similar to enforcing a ban on proprietary data and algorithms. One example is Article 20 of the General Data Protection Regulation \parencite{de2018right}.} or a floor on the monitoring discount---would further reduce monitoring take-up. This ends up hurting consumers by weakening the firm's dynamic incentives to invest and by raising reclassification risk in equilibrium.

The paper proceeds as follows. Section 1 describes our data and provides background information on auto insurance and monitoring. Section 2 conducts reduced-form tests to evaluate the effects of moral hazard and selection that are induced by the monitoring program. Section 3 presents our structural model, identification arguments, and estimation procedures to recover key cost and demand parameters. Section 4 discusses estimation results. Section 5 proposes a model of monitoring pricing and studies the welfare implications of the monitoring program in the status quo as well as counterfactuals involving optimal pricing and a ban on proprietary data. Section 6 revisits our results in relation to the literature and concludes.
